
\documentclass[aspectratio=169]{beamer}
\usetheme{Madrid}
\usecolortheme{default}
\usefonttheme{professionalfonts}
\setbeamertemplate{navigation symbols}{}
\setbeamertemplate{footline}[frame number]
\setbeamertemplate{caption}[numbered]

% --- Packages ---
\usepackage[utf8]{inputenc}
\usepackage[T1]{fontenc}
\usepackage[italian]{babel}
\usepackage{siunitx}
\usepackage{amsmath, amsfonts, amssymb}
\usepackage{physics}
\usepackage{graphicx}
\usepackage{booktabs}
\usepackage{hyperref}
\hypersetup{colorlinks=true, urlcolor=blue, linkcolor=blue}

% --- Title ---
\title{Derivatore e Integratore (S05)}
\subtitle{Analisi teorica, simulazioni e risultati sperimentali}
\author{TD - Tecnologie Digitali}
\institute{OpAmp MCP601 -- Laboratorio}
\date{\today}

\begin{document}

% --- Title page ---
\begin{frame}
  \titlepage
\end{frame}

% --- Outline ---
\begin{frame}{Outline}
\tableofcontents
\end{frame}

\section{Obiettivi e schema dell'esercitazione}
\begin{frame}{Obiettivi S05}
\begin{itemize}
  \item Ricavare e discutere la funzione di trasferimento $H(\omega)$ di \textbf{integratore} e \textbf{derivatore}.
  \item Stimare il \textbf{range di frequenze} in cui i circuiti operano come integratore/derivatore ideali.
  \item Valutare l'effetto dell'\textbf{offset} $V_{\text{off}}$ sull'uscita e il rischio di \textbf{saturazione}.
  \item Misurare la \textbf{risposta in frequenza} e confrontarla con la teoria (diagrammi di Bode).
\end{itemize}
\end{frame}

\begin{frame}{Scheda: integratore e derivatore di Miller}
\begin{columns}[T,onlytextwidth]
\column{0.5\textwidth}
\begin{itemize}
  \item Integratore: rete $RC$ in \textit{feedback}.
  \item Derivatore: rete $RC$ in \textit{ingresso}.
  \item Resistenze di regolazione: $R_{\text{reg}}$ (stabilizza, limita guadagno DC).
\end{itemize}
\column{0.5\textwidth}
\centering
\includegraphics[width=\linewidth]{/mnt/data/c0aebc44-2399-44ae-a981-9d6a8b26264f.png}
\end{columns}
\end{frame}

\section{Teoria essenziale}
\begin{frame}{Integratore ideale e con $R_{\text{reg}}$}
\textbf{Integratore ideale} (ipotesi $A\to\infty$ e ingresso dell'OpAmp a massa virtuale):
\begin{equation}
H(\omega)=\frac{V_{\text{out}}}{V_{\text{in}}} = -\frac{1}{j\omega R C}.
\end{equation}
Inserendo la resistenza di feedback $R_{\text{reg}}$ si ottiene un polo a bassa frequenza e un guadagno finito a DC:
\begin{equation}
H(\omega)=\frac{R_{\text{reg}}/R}{1 + j\omega R_{\text{reg}} C}, \qquad
\omega \gg \frac{1}{R_{\text{reg}} C} \Rightarrow \text{comportamento da integratore}.
\end{equation}
Nel dominio del tempo: \quad
$V_{\text{out}}(t) = -\dfrac{1}{RC}\displaystyle\int V_{\text{in}}(t)\,dt$ (trascurando transitori).
\end{frame}

\begin{frame}{Derivatore con polo dominante dell'amplificatore}
Derivatore ideale: 
\begin{equation}
H(\omega) = - j \omega R C.
\end{equation}
Considerando la \textbf{dipendenza in frequenza del guadagno open-loop} dell'OpAmp
\begin{equation}
A(\omega) = \frac{A_0}{1 + j\omega\tau_{PD}}, \qquad \text{(polo dominante)}
\end{equation}
si ha, per $\abs{A(\omega)}\gg 1$,
\begin{equation}
H(\omega) \simeq \frac{A_0\,j\omega\tau_{PD}}{1 + j\omega\tau_{PD}} \cdot \frac{R}{A_0} = \frac{j\omega R C}{1 + j\omega\tau_{PD}} \quad \text{(forma qualitativa)}.
\end{equation}
Ne segue una \textbf{divergenza limitata} e un \textbf{massimo} attorno a
$f_{\text{div}} \propto \dfrac{1}{2\pi\sqrt{\tau_{PD} R C A_0}}$.
\end{frame}

\section{Integratore: simulazioni e prove}
\begin{frame}{Montaggio integratore e simulazioni (TINA)}
\begin{columns}[T,onlytextwidth]
\column{0.5\textwidth}
\includegraphics[width=\linewidth]{/mnt/data/5ece17af-bb8c-448c-bed8-77bdf3c26beb.png}
\column{0.5\textwidth}
\begin{itemize}
  \item Simulazioni di spazzata in frequenza e nel tempo.
  \item Verifica con onda \textbf{quadra} $\rightarrow$ uscita \textbf{triangolare}.
  \item Valutazione del limite a bassa $f$ fissato da $R_{\text{reg}}$.
\end{itemize}
\end{columns}
\end{frame}

\begin{frame}{Integratore: esempi dal logbook (tempo)}
\centering
\includegraphics[page=8,width=.9\linewidth]{/mnt/data/MALDIN_logbook_es05.pdf}\\[2mm]
\small Esempi ingresso/uscita a diverse frequenze (uscita triangolare per ingresso quadra).
\end{frame}

\begin{frame}{Integratore: risposta in frequenza}
\centering
\includegraphics[page=20,width=.9\linewidth]{/mnt/data/MALDIN_logbook_es05.pdf}\\[2mm]
\small Spazzata in frequenza: regione di comportamento da integratore tra $\frac{1}{2\pi R_{\text{reg}}C}$ e la banda utile dell'OpAmp.
\end{frame}

\begin{frame}{Offset $V_{\text{off}}$ e saturazione (integratore)}
Nel limite di alta amplificazione, l'offset d'ingresso si ripresenta all'ingresso invertente:
\begin{equation}
V_{+}=V_{-}=V_{\text{off}}.
\end{equation}
La componente continua in uscita cresce come
\begin{equation}
V_{\text{out,DC}} \simeq V_{\text{off}}\,\frac{R+R_{\text{reg}}}{R},
\end{equation}
portando a \textbf{saturazione} se non si limita l'intervallo di $f$ o l'ampiezza del segnale.
\end{frame}

\section{Derivatore: simulazioni e prove}
\begin{frame}{Montaggio derivatore e stabilizzazione}
\begin{columns}[T,onlytextwidth]
\column{0.5\textwidth}
\includegraphics[width=\linewidth]{/mnt/data/2d0494ac-467e-4326-ab4a-cf51b54fcbb0.png}
\column{0.5\textwidth}
\begin{itemize}
  \item $R_{\text{reg}}$ in serie/parallelo per limitare guadagno in alta $f$.
  \item Evita auto-oscillazioni e risonanza marcata.
\end{itemize}
\end{columns}
\end{frame}

\begin{frame}{Derivatore: esempi dal logbook (tempo)}
\centering
\includegraphics[page=15,width=.9\linewidth]{/mnt/data/MALDIN_logbook_es05.pdf}\\[2mm]
\small Uscita proporzionale alla derivata dell'ingresso: picchi in corrispondenza dei fronti.
\end{frame}

\begin{frame}{Derivatore: risonanza e mitigazione}
\centering
\includegraphics[page=17,width=.9\linewidth]{/mnt/data/MALDIN_logbook_es05.pdf}\\[2mm]
\small Campionamento della risonanza senza $R_{\text{reg}}$; l'inserimento di $R_{\text{reg}}$ smorza e allarga la banda utile.
\end{frame}

\section{Confronti in frequenza}
\begin{frame}{Bode: integrazione vs derivazione}
\begin{columns}[T,onlytextwidth]
\column{0.5\textwidth}
\centering
\includegraphics[page=20,width=\linewidth]{/mnt/data/MALDIN_logbook_es05.pdf}
\small Integratore: pendenza $\approx -20$ dB/dec oltre $\frac{1}{R_{\text{reg}}C}$.
\column{0.5\textwidth}
\centering
\includegraphics[page=21,width=\linewidth]{/mnt/data/MALDIN_logbook_es05.pdf}
\small Derivatore: pendenza $\approx +20$ dB/dec fino al limite imposto da $A(\omega)$ e $\tau_{PD}$.
\end{columns}
\end{frame}

\section{Takeaways e prospettive}
\begin{frame}{Cosa abbiamo imparato}
\begin{itemize}
  \item \textbf{Integratore}: funziona per $\omega \gg 1/(R_{\text{reg}}C)$, ma è sensibile all'offset $\Rightarrow$ rischio saturazione.
  \item \textbf{Derivatore}: amplifica il rumore alle alte frequenze; $R_{\text{reg}}$ è \textbf{necessaria} per stabilità.
  \item Il polo dominante dell'OpAmp limita entrambi i circuiti: attenzione a \textbf{GBWP} e a $\tau_{PD}$.
  \item Le misure confermano le tendenze teoriche con \textbf{deviazioni spiegabili} da guadagno finito e componenti reali.
\end{itemize}
\end{frame}

\begin{frame}{Appendice: note pratiche}
\begin{itemize}
  \item Evitare saturazione regolando ampiezza del generatore e scegliendo $R_{\text{reg}}$ adeguata.
  \item Per misure di Bode pulite: bufferare l'ingresso e curare la messa a terra.
  \item Consiglio: salvare tracce e parametri di setup per tracciabilità.
\end{itemize}
\end{frame}

\begin{frame}
\centering
\Large Domande?
\end{frame}

\end{document}
"""

tex_path.write_text(content, encoding="utf-8")

# Also create a simple README to guide compilation
readme = base / "README_S05.txt"
readme.write_text("""
Compilazione:
1) Usare pdflatex o lualatex su S05_derivatore_integratore_beamer.tex
2) Il file include direttamente alcune pagine del PDF del logbook e 4 immagini PNG fornite.
   - Se lavori in un'altra cartella, mantieni il PDF 'MALDIN_logbook_es05.pdf' nella stessa directory
     oppure aggiorna i percorsi nelle \\includegraphics.
""", encoding="utf-8")

[str(tex_path), str(readme)]
